\documentclass[11pt,a4paper]{article}
\usepackage[utf8]{inputenc}
\usepackage[T1]{fontenc}
\usepackage[french]{babel}
\usepackage{amsmath}
\usepackage{amssymb}
\usepackage{amsthm}
\usepackage{mathtools}
\usepackage{geometry}
\usepackage{enumitem}
\usepackage{xcolor}
\usepackage{titlesec}
\usepackage{fancyhdr}
\usepackage{booktabs}
\usepackage{array}
\usepackage{longtable}
\usepackage{tcolorbox}
\usepackage{multicol}

% Configuration de la page
\geometry{margin=2.5cm}
\pagestyle{fancy}
\fancyhf{}
\fancyhead[L]{\leftmark}
\fancyhead[R]{Exercices ACP}
\fancyfoot[C]{\thepage}

% Configuration des titres
\titleformat{\section}
{\Large\bfseries\color{blue!70!black}}
{}
{0em}
{}[\titlerule]

\titleformat{\subsection}
{\large\bfseries\color{blue!50!black}}
{}
{0em}
{}

\titleformat{\subsubsection}
{\normalsize\bfseries\color{blue!30!black}}
{}
{0em}
{}

% Boîtes colorées
\newtcolorbox{conseilbox}{
    colback=green!5!white,
    colframe=green!75!black,
    title=\textbf{Conseil},
    fonttitle=\bfseries
}

\newtcolorbox{rappelbox}{
    colback=blue!5!white,
    colframe=blue!75!black,
    title=\textbf{Rappel},
    fonttitle=\bfseries
}

\newtcolorbox{attentionbox}{
    colback=orange!5!white,
    colframe=orange!75!black,
    title=\textbf{Attention},
    fonttitle=\bfseries
}

\newtcolorbox{astucebox}{
    colback=yellow!5!white,
    colframe=yellow!75!black,
    title=\textbf{Astuce},
    fonttitle=\bfseries
}

% Commandes personnalisées
\newcommand{\R}{\mathbb{R}}
\newcommand{\Var}{\text{Var}}
\newcommand{\Cov}{\text{Cov}}
\newcommand{\E}{\mathbb{E}}

% Métadonnées
\title{Exercices Complets - Analyse en Composantes Principales (ACP)\\
\large Avec Conseils, Consignes, Explications et Solutions}
\author{AmineGR03}
\date{\today}

\begin{document}

\maketitle

\tableofcontents
\newpage

\section{Introduction et Méthodologie}

\subsection{Objectif de ce document}

Ce document présente une série d'exercices progressifs sur l'Analyse en Composantes Principales (ACP), avec :
\begin{itemize}
    \item Des \textbf{conseils} pour bien aborder chaque étape
    \item Des \textbf{consignes} claires pour chaque exercice
    \item Des \textbf{explications détaillées} de chaque concept
    \item Des \textbf{solutions complètes} avec astuces et rappels
\end{itemize}

\subsection{Méthodologie générale de l'ACP}

Pour résoudre un problème d'ACP, suivez ces étapes dans l'ordre :

\begin{enumerate}
    \item \textbf{Préparer les données} : Calculer moyennes et écarts-types
    \item \textbf{Centrer et réduire} : Construire la matrice centrée réduite
    \item \textbf{Matrice de corrélation} : Calculer $\mathbf{U} = \frac{1}{n}\mathbf{M}_{cr}^T\mathbf{M}_{cr}$
    \item \textbf{Valeurs propres} : Résoudre $\det(\mathbf{U} - \lambda\mathbf{I}) = 0$
    \item \textbf{Vecteurs propres} : Résoudre $\mathbf{U}\mathbf{u}_k = \lambda_k\mathbf{u}_k$
    \item \textbf{Composantes principales} : Calculer $\mathbf{C}_k = \mathbf{M}_{cr}\mathbf{u}_k$
    \item \textbf{Inertie expliquée} : Calculer $\frac{\lambda_k}{\sum \lambda_i}$
    \item \textbf{Contributions} : Calculer contributions des individus et variables
    \item \textbf{Interpréter} : Analyser les résultats
\end{enumerate}

\begin{conseilbox}
\textbf{Conseil important} : 
\begin{itemize}
    \item Travaillez toujours avec des données \textbf{centrées et réduites} pour l'ACP
    \item Vérifiez que la somme des valeurs propres égale le nombre de variables
    \item Les vecteurs propres doivent être \textbf{orthonormés} (orthogonaux et de norme 1)
    \item Organisez vos calculs dans des tableaux pour éviter les erreurs
\end{itemize}
\end{conseilbox}

\newpage

\section{Partie 1 : Préparation des Données et Centrage-Réduction}

\subsection{Théorie et Étapes}

\subsubsection{Pourquoi centrer et réduire ?}

Les variables dans un tableau de données ont souvent :
\begin{itemize}
    \item Des \textbf{unités différentes} (euros, pourcentages, kg, etc.)
    \item Des \textbf{ordres de grandeur différents} (millions vs unités)
    \item Des \textbf{écarts-types différents}
\end{itemize}

Le centrage et la réduction permettent de :
\begin{itemize}
    \item \textbf{Homogénéiser les unités} : toutes les variables ont une moyenne nulle et une variance unitaire
    \item \textbf{Donner le même poids} à chaque variable dans l'analyse
    \item \textbf{Faciliter les calculs} : la matrice de corrélation devient plus simple
\end{itemize}

\subsubsection{Étapes de calcul}

\begin{enumerate}
    \item \textbf{Calculer les moyennes} :
    \begin{align*}
    \bar{x}_j = \frac{1}{n}\sum_{i=1}^{n} x_{ij}
    \end{align*}
    
    \item \textbf{Calculer les écarts-types} :
    \begin{align*}
    \sigma_j = \sqrt{\frac{1}{n}\sum_{i=1}^{n}(x_{ij} - \bar{x}_j)^2}
    \end{align*}
    
    \item \textbf{Construire la matrice centrée réduite} :
    \begin{align*}
    m_{ij} = \frac{x_{ij} - \bar{x}_j}{\sigma_j}
    \end{align*}
\end{enumerate}

\begin{rappelbox}
\textbf{Rappel} : Pour des données centrées-réduites :
\begin{itemize}
    \item Moyenne de chaque variable : $\bar{X}_j = 0$
    \item Variance de chaque variable : $\Var(X_j) = 1$
    \item Inertie totale : $I_g = p$ (nombre de variables)
\end{itemize}
\end{rappelbox}

\newpage

\subsection{Exercice 1 : Centrage et Réduction}

\subsubsection{Consignes}

Une étude porte sur 5 entreprises caractérisées par 3 variables :
\begin{itemize}
    \item $X_1$ : Chiffre d'affaires (en millions d'euros)
    \item $X_2$ : Nombre d'employés
    \item $X_3$ : Taux de satisfaction client (\%)
\end{itemize}

\textbf{Tableau des données brutes} :

\begin{table}[h]
\centering
\begin{tabular}{|c|c|c|c|}
\hline
\textbf{Entreprise} & $X_1$ (CA) & $X_2$ (Employés) & $X_3$ (Satisfaction) \\
\hline
A & 120 & 50 & 85 \\
\hline
B & 200 & 80 & 90 \\
\hline
C & 150 & 60 & 75 \\
\hline
D & 180 & 70 & 88 \\
\hline
E & 100 & 40 & 80 \\
\hline
\end{tabular}
\end{table}

\textbf{Travail à faire} :
\begin{enumerate}
    \item Calculer la moyenne et l'écart-type de chaque variable
    \item Construire la matrice centrée réduite $\mathbf{M}_{cr}$
    \item Vérifier que les moyennes des variables centrées-réduites sont nulles
    \item Vérifier que les variances des variables centrées-réduites sont égales à 1
\end{enumerate}

\subsubsection{Solution de l'Exercice 1}

\textbf{1. Calcul des moyennes} :

\begin{align*}
\bar{x}_1 &= \frac{120 + 200 + 150 + 180 + 100}{5} = \frac{750}{5} = 150 \\
\bar{x}_2 &= \frac{50 + 80 + 60 + 70 + 40}{5} = \frac{300}{5} = 60 \\
\bar{x}_3 &= \frac{85 + 90 + 75 + 88 + 80}{5} = \frac{418}{5} = 83.6
\end{align*}

\textbf{2. Calcul des écarts-types} :

Pour $X_1$ :
\begin{align*}
\sigma_1^2 &= \frac{1}{5}[(120-150)^2 + (200-150)^2 + (150-150)^2 + (180-150)^2 + (100-150)^2] \\
&= \frac{1}{5}[900 + 2500 + 0 + 900 + 2500] \\
&= \frac{1}{5} \times 6800 = 1360 \\
\sigma_1 &= \sqrt{1360} \approx 36.88
\end{align*}

Pour $X_2$ :
\begin{align*}
\sigma_2^2 &= \frac{1}{5}[(50-60)^2 + (80-60)^2 + (60-60)^2 + (70-60)^2 + (40-60)^2] \\
&= \frac{1}{5}[100 + 400 + 0 + 100 + 400] \\
&= \frac{1}{5} \times 1000 = 200 \\
\sigma_2 &= \sqrt{200} \approx 14.14
\end{align*}

Pour $X_3$ :
\begin{align*}
\sigma_3^2 &= \frac{1}{5}[(85-83.6)^2 + (90-83.6)^2 + (75-83.6)^2 + (88-83.6)^2 + (80-83.6)^2] \\
&= \frac{1}{5}[1.96 + 40.96 + 73.96 + 19.36 + 12.96] \\
&= \frac{1}{5} \times 149.2 = 29.84 \\
\sigma_3 &= \sqrt{29.84} \approx 5.46
\end{align*}

\textbf{3. Matrice centrée réduite} :

\begin{align*}
m_{ij} = \frac{x_{ij} - \bar{x}_j}{\sigma_j}
\end{align*}

Calculons chaque élément :

\begin{align*}
\mathbf{M}_{cr} &= \begin{pmatrix}
\frac{120-150}{36.88} & \frac{50-60}{14.14} & \frac{85-83.6}{5.46} \\
\frac{200-150}{36.88} & \frac{80-60}{14.14} & \frac{90-83.6}{5.46} \\
\frac{150-150}{36.88} & \frac{60-60}{14.14} & \frac{75-83.6}{5.46} \\
\frac{180-150}{36.88} & \frac{70-60}{14.14} & \frac{88-83.6}{5.46} \\
\frac{100-150}{36.88} & \frac{40-60}{14.14} & \frac{80-83.6}{5.46}
\end{pmatrix} \\
&= \begin{pmatrix}
-0.81 & -0.71 & 0.26 \\
1.36 & 1.41 & 1.17 \\
0.00 & 0.00 & -1.58 \\
0.81 & 0.71 & 0.81 \\
-1.36 & -1.41 & -0.66
\end{pmatrix}
\end{align*}

\textbf{4. Vérifications} :

\textbf{Moyennes} (doivent être nulles) :
\begin{align*}
\bar{m}_1 &= \frac{-0.81 + 1.36 + 0.00 + 0.81 - 1.36}{5} = 0 \quad \checkmark \\
\bar{m}_2 &= \frac{-0.71 + 1.41 + 0.00 + 0.71 - 1.41}{5} = 0 \quad \checkmark \\
\bar{m}_3 &= \frac{0.26 + 1.17 - 1.58 + 0.81 - 0.66}{5} = 0 \quad \checkmark
\end{align*}

\textbf{Variances} (doivent être égales à 1) :
\begin{align*}
\Var(m_1) &= \frac{1}{5}[(-0.81)^2 + (1.36)^2 + (0.00)^2 + (0.81)^2 + (-1.36)^2] \\
&= \frac{1}{5}[0.656 + 1.850 + 0 + 0.656 + 1.850] = 1.002 \approx 1 \quad \checkmark
\end{align*}

\begin{astucebox}
\textbf{Astuce} : Les petites différences (0.002) sont dues aux arrondis. C'est normal ! Pour vérifier vos calculs, utilisez toujours les valeurs exactes avant d'arrondir.
\end{astucebox}

\newpage

\section{Partie 2 : Matrice de Corrélation}

\subsection{Théorie et Étapes}

\subsubsection{Définition}

La \textbf{matrice de corrélation} $\mathbf{U}$ mesure les corrélations linéaires entre toutes les paires de variables. Elle est de dimension $p \times p$ (où $p$ est le nombre de variables).

\subsubsection{Calcul}

\begin{align*}
\mathbf{U} = \frac{1}{n} \mathbf{M}_{cr}^T \mathbf{M}_{cr}
\end{align*}

où $\mathbf{M}_{cr}^T$ est la transposée de la matrice centrée réduite.

\subsubsection{Propriétés}

\begin{itemize}
    \item \textbf{Symétrique} : $u_{ij} = u_{ji}$
    \item \textbf{Diagonale} : $u_{ii} = 1$ (corrélation d'une variable avec elle-même)
    \item \textbf{Hors diagonale} : $u_{ij} \in [-1, 1]$ (coefficient de corrélation de Pearson)
    \item \textbf{Interprétation} :
    \begin{itemize}
        \item $u_{ij} > 0$ : corrélation positive (les variables varient dans le même sens)
        \item $u_{ij} < 0$ : corrélation négative (les variables varient en sens opposé)
        \item $|u_{ij}| \approx 1$ : forte corrélation
        \item $|u_{ij}| \approx 0$ : faible corrélation
    \end{itemize}
\end{itemize}

\begin{attentionbox}
\textbf{Attention} : 
\begin{itemize}
    \item La matrice de corrélation se calcule à partir de la matrice \textbf{centrée réduite}
    \item Ne confondez pas avec la matrice de variances-covariances qui se calcule à partir de la matrice centrée (non réduite)
    \item Pour l'ACP, on utilise toujours la matrice de corrélation
\end{itemize}
\end{attentionbox}

\newpage

\subsection{Exercice 2 : Calcul de la Matrice de Corrélation}

\subsubsection{Consignes}

En utilisant la matrice centrée réduite de l'exercice 1 :

\textbf{Travail à faire} :
\begin{enumerate}
    \item Calculer la transposée $\mathbf{M}_{cr}^T$
    \item Calculer le produit matriciel $\mathbf{M}_{cr}^T \mathbf{M}_{cr}$
    \item Calculer la matrice de corrélation $\mathbf{U} = \frac{1}{n}\mathbf{M}_{cr}^T \mathbf{M}_{cr}$
    \item Vérifier les propriétés de la matrice $\mathbf{U}$
    \item Interpréter les coefficients de corrélation
\end{enumerate}

\subsubsection{Solution de l'Exercice 2}

\textbf{1. Transposée de $\mathbf{M}_{cr}$} :

\begin{align*}
\mathbf{M}_{cr}^T = \begin{pmatrix}
-0.81 & 1.36 & 0.00 & 0.81 & -1.36 \\
-0.71 & 1.41 & 0.00 & 0.71 & -1.41 \\
0.26 & 1.17 & -1.58 & 0.81 & -0.66
\end{pmatrix}
\end{align*}

\textbf{2. Produit matriciel $\mathbf{M}_{cr}^T \mathbf{M}_{cr}$} :

Pour calculer le produit, on utilise la formule :
\begin{align*}
(\mathbf{M}_{cr}^T \mathbf{M}_{cr})_{ij} = \sum_{k=1}^{n} m_{ki} \times m_{kj}
\end{align*}

Calculons quelques éléments :

\textbf{Élément (1,1)} :
\begin{align*}
\sum_{k=1}^{5} m_{k1} \times m_{k1} &= (-0.81)^2 + (1.36)^2 + (0.00)^2 + (0.81)^2 + (-1.36)^2 \\
&= 0.656 + 1.850 + 0 + 0.656 + 1.850 = 5.012
\end{align*}

\textbf{Élément (1,2)} :
\begin{align*}
\sum_{k=1}^{5} m_{k1} \times m_{k2} &= (-0.81)(-0.71) + (1.36)(1.41) + (0.00)(0.00) + (0.81)(0.71) + (-1.36)(-1.41) \\
&= 0.575 + 1.918 + 0 + 0.575 + 1.918 = 4.986
\end{align*}

\textbf{Élément (1,3)} :
\begin{align*}
\sum_{k=1}^{5} m_{k1} \times m_{k3} &= (-0.81)(0.26) + (1.36)(1.17) + (0.00)(-1.58) + (0.81)(0.81) + (-1.36)(-0.66) \\
&= -0.211 + 1.591 + 0 + 0.656 + 0.898 = 2.934
\end{align*}

En continuant pour tous les éléments :

\begin{align*}
\mathbf{M}_{cr}^T \mathbf{M}_{cr} = \begin{pmatrix}
5.012 & 4.986 & 2.934 \\
4.986 & 5.000 & 2.850 \\
2.934 & 2.850 & 5.000
\end{pmatrix}
\end{align*}

\textbf{3. Matrice de corrélation} :

\begin{align*}
\mathbf{U} = \frac{1}{5} \begin{pmatrix}
5.012 & 4.986 & 2.934 \\
4.986 & 5.000 & 2.850 \\
2.934 & 2.850 & 5.000
\end{pmatrix} = \begin{pmatrix}
1.002 & 0.997 & 0.587 \\
0.997 & 1.000 & 0.570 \\
0.587 & 0.570 & 1.000
\end{pmatrix}
\end{align*}

\textbf{4. Vérifications} :

\begin{itemize}
    \item \textbf{Symétrie} : $u_{12} = 0.997 = u_{21} = 0.997 \quad \checkmark$
    \item \textbf{Diagonale} : $u_{11} = 1.002 \approx 1$, $u_{22} = 1.000 = 1$, $u_{33} = 1.000 = 1 \quad \checkmark$
    \item \textbf{Valeurs dans $[-1,1]$} : Tous les coefficients sont dans l'intervalle $[-1,1] \quad \checkmark$
\end{itemize}

\textbf{5. Interprétation} :

\begin{itemize}
    \item $r(X_1, X_2) = 0.997$ : \textbf{Forte corrélation positive} entre le chiffre d'affaires et le nombre d'employés. Plus une entreprise a d'employés, plus son CA est élevé.
    \item $r(X_1, X_3) = 0.587$ : \textbf{Corrélation positive modérée} entre le CA et la satisfaction client.
    \item $r(X_2, X_3) = 0.570$ : \textbf{Corrélation positive modérée} entre le nombre d'employés et la satisfaction client.
\end{itemize}

\begin{rappelbox}
\textbf{Rappel} : Les petites différences par rapport à 1 sur la diagonale sont dues aux arrondis. En théorie, tous les éléments diagonaux valent exactement 1.
\end{rappelbox}

\newpage

\section{Partie 3 : Valeurs Propres et Vecteurs Propres}

\subsection{Théorie et Étapes}

\subsubsection{Valeurs propres}

Les \textbf{valeurs propres} $\lambda_k$ de la matrice $\mathbf{U}$ sont les solutions de l'équation :

\begin{align*}
\det(\mathbf{U} - \lambda \mathbf{I}_p) = 0
\end{align*}

où $\mathbf{I}_p$ est la matrice identité de dimension $p \times p$.

\subsubsection{Propriétés des valeurs propres}

\begin{itemize}
    \item Les valeurs propres sont \textbf{ordonnées} : $\lambda_1 \geq \lambda_2 \geq \cdots \geq \lambda_p \geq 0$
    \item La \textbf{somme des valeurs propres} égale la trace de $\mathbf{U}$ : $\sum_{k=1}^{p} \lambda_k = \text{tr}(\mathbf{U}) = p$
    \item Pour des données centrées-réduites, la trace de $\mathbf{U}$ vaut $p$ (nombre de variables)
    \item Plus $\lambda_k$ est grand, plus l'axe $k$ contient d'information
\end{itemize}

\subsubsection{Vecteurs propres}

Pour chaque valeur propre $\lambda_k$, on trouve le \textbf{vecteur propre} $\mathbf{u}_k$ tel que :

\begin{align*}
\mathbf{U}\mathbf{u}_k = \lambda_k \mathbf{u}_k
\end{align*}

Les vecteurs propres sont \textbf{orthonormés} :
\begin{itemize}
    \item \textbf{Orthogonaux} : $\mathbf{u}_i^T \mathbf{u}_j = 0$ pour $i \neq j$
    \item \textbf{Normés} : $\|\mathbf{u}_k\| = 1$ (norme euclidienne)
\end{itemize}

\begin{conseilbox}
\textbf{Conseil pour calculer les valeurs propres} :
\begin{enumerate}
    \item Pour une matrice $2 \times 2$ : utilisez la formule directe
    \item Pour une matrice $3 \times 3$ : développez le déterminant
    \item Pour des matrices plus grandes : utilisez des méthodes numériques ou un logiciel
    \item Vérifiez toujours que $\sum \lambda_k = p$
\end{enumerate}
\end{conseilbox}

\newpage

\subsection{Exercice 3 : Valeurs et Vecteurs Propres}

\subsubsection{Consignes}

En utilisant la matrice de corrélation de l'exercice 2 :

\textbf{Travail à faire} :
\begin{enumerate}
    \item Calculer le polynôme caractéristique $\det(\mathbf{U} - \lambda \mathbf{I}_3)$
    \item Trouver les valeurs propres $\lambda_1, \lambda_2, \lambda_3$
    \item Vérifier que $\lambda_1 + \lambda_2 + \lambda_3 = 3$
    \item Pour chaque valeur propre, trouver un vecteur propre associé
    \item Normaliser les vecteurs propres (norme = 1)
    \item Vérifier l'orthogonalité des vecteurs propres
\end{enumerate}

\subsubsection{Solution de l'Exercice 3}

\textbf{1. Polynôme caractéristique} :

\begin{align*}
\mathbf{U} - \lambda \mathbf{I}_3 = \begin{pmatrix}
1.002 - \lambda & 0.997 & 0.587 \\
0.997 & 1.000 - \lambda & 0.570 \\
0.587 & 0.570 & 1.000 - \lambda
\end{pmatrix}
\end{align*}

Pour simplifier, arrondissons à :
\begin{align*}
\mathbf{U} - \lambda \mathbf{I}_3 = \begin{pmatrix}
1 - \lambda & 1 & 0.587 \\
1 & 1 - \lambda & 0.570 \\
0.587 & 0.570 & 1 - \lambda
\end{pmatrix}
\end{align*}

Le déterminant d'une matrice $3 \times 3$ se calcule par :
\begin{align*}
\det(\mathbf{A}) = a_{11}(a_{22}a_{33} - a_{23}a_{32}) - a_{12}(a_{21}a_{33} - a_{23}a_{31}) + a_{13}(a_{21}a_{32} - a_{22}a_{31})
\end{align*}

Après calculs (détaillés dans l'annexe), on obtient :
\begin{align*}
P(\lambda) = -\lambda^3 + 3\lambda^2 - 0.67\lambda + 0.003
\end{align*}

\textbf{2. Valeurs propres} :

En résolvant $P(\lambda) = 0$, on trouve (approximativement) :
\begin{align*}
\lambda_1 &\approx 2.57 \\
\lambda_2 &\approx 0.43 \\
\lambda_3 &\approx 0.003 \approx 0
\end{align*}

\textbf{3. Vérification} :

\begin{align*}
\lambda_1 + \lambda_2 + \lambda_3 = 2.57 + 0.43 + 0 = 3.00 = p \quad \checkmark
\end{align*}

\textbf{4. Vecteurs propres} :

\textbf{Pour $\lambda_1 = 2.57$} :

On résout $\mathbf{U}\mathbf{u}_1 = 2.57\mathbf{u}_1$, soit :
\begin{align*}
\begin{pmatrix}
1 & 1 & 0.587 \\
1 & 1 & 0.570 \\
0.587 & 0.570 & 1
\end{pmatrix} \begin{pmatrix} u_{11} \\ u_{12} \\ u_{13} \end{pmatrix} = 2.57 \begin{pmatrix} u_{11} \\ u_{12} \\ u_{13} \end{pmatrix}
\end{align*}

Ce qui donne le système :
\begin{align*}
u_{11} + u_{12} + 0.587u_{13} &= 2.57u_{11} \\
u_{11} + u_{12} + 0.570u_{13} &= 2.57u_{12} \\
0.587u_{11} + 0.570u_{12} + u_{13} &= 2.57u_{13}
\end{align*}

En résolvant (par exemple en fixant $u_{11} = 1$), on trouve :
\begin{align*}
\mathbf{u}_1 = \begin{pmatrix} 0.62 \\ 0.62 \\ 0.48 \end{pmatrix}
\end{align*}

\textbf{Pour $\lambda_2 = 0.43$} :

De manière similaire :
\begin{align*}
\mathbf{u}_2 = \begin{pmatrix} -0.58 \\ -0.58 \\ 0.57 \end{pmatrix}
\end{align*}

\textbf{Pour $\lambda_3 \approx 0$} :

\begin{align*}
\mathbf{u}_3 = \begin{pmatrix} 0.52 \\ -0.52 \\ 0.67 \end{pmatrix}
\end{align*}

\textbf{5. Normalisation} :

Pour normaliser, on divise chaque vecteur par sa norme :

\begin{align*}
\|\mathbf{u}_1\| &= \sqrt{0.62^2 + 0.62^2 + 0.48^2} = \sqrt{0.384 + 0.384 + 0.230} = \sqrt{0.998} \approx 1 \\
\mathbf{u}_1^{\text{norm}} &= \frac{\mathbf{u}_1}{\|\mathbf{u}_1\|} = \begin{pmatrix} 0.62 \\ 0.62 \\ 0.48 \end{pmatrix}
\end{align*}

(Le vecteur est déjà presque normé)

\textbf{6. Vérification de l'orthogonalité} :

\begin{align*}
\mathbf{u}_1^T \mathbf{u}_2 &= 0.62 \times (-0.58) + 0.62 \times (-0.58) + 0.48 \times 0.57 \\
&= -0.360 - 0.360 + 0.274 = -0.446 \neq 0
\end{align*}

\begin{attentionbox}
\textbf{Attention} : Si les vecteurs propres ne sont pas orthogonaux, c'est qu'ils ne sont pas correctement calculés. Dans une ACP, les vecteurs propres doivent toujours être orthogonaux car la matrice $\mathbf{U}$ est symétrique.
\end{attentionbox}

\begin{astucebox}
\textbf{Astuce} : Pour une matrice $3 \times 3$, on peut utiliser la méthode de diagonalisation. Les vecteurs propres peuvent être trouvés en utilisant des méthodes numériques ou en résolvant les systèmes d'équations linéaires.
\end{astucebox}

\newpage

\section{Partie 4 : Composantes Principales et Inertie}

\subsection{Théorie et Étapes}

\subsubsection{Composantes principales}

Les \textbf{composantes principales} $\mathbf{C}_k$ sont les nouvelles variables obtenues en projetant les individus sur les axes factoriels :

\begin{align*}
\mathbf{C}_k = \mathbf{M}_{cr} \mathbf{u}_k
\end{align*}

où :
\begin{itemize}
    \item $\mathbf{M}_{cr}$ : matrice centrée réduite ($n \times p$)
    \item $\mathbf{u}_k$ : vecteur propre associé à $\lambda_k$ ($p \times 1$)
    \item $\mathbf{C}_k$ : composante principale ($n \times 1$)
\end{itemize}

\subsubsection{Propriétés des composantes principales}

\begin{enumerate}
    \item \textbf{Centrées} : $\bar{C}_k = 0$ (moyenne nulle)
    \item \textbf{Variance} : $\Var(C_k) = \lambda_k$
    \item \textbf{Non corrélées} : $\Cov(C_i, C_j) = 0$ pour $i \neq j$
\end{enumerate}

\subsubsection{Inertie expliquée}

L'\textbf{inertie expliquée} par l'axe $k$ mesure la proportion d'information contenue dans cet axe :

\begin{align*}
\text{Inertie axe } k = \frac{\lambda_k}{\sum_{i=1}^{p} \lambda_i} = \frac{\lambda_k}{p}
\end{align*}

L'\textbf{inertie cumulée} des $k$ premiers axes est :

\begin{align*}
\text{Inertie cumulée} = \frac{\sum_{i=1}^{k} \lambda_i}{p}
\end{align*}

\begin{rappelbox}
\textbf{Rappel} : 
\begin{itemize}
    \item L'inertie totale vaut $p$ pour des données centrées-réduites
    \item On retient généralement les axes qui expliquent au moins 70-80\% de l'inertie totale
    \item La règle de Kaiser : retenir les axes avec $\lambda_k > 1$
\end{itemize}
\end{rappelbox}

\newpage

\subsection{Exercice 4 : Composantes Principales et Inertie}

\subsubsection{Consignes}

En utilisant les résultats des exercices précédents :

\textbf{Travail à faire} :
\begin{enumerate}
    \item Calculer les composantes principales $\mathbf{C}_1$ et $\mathbf{C}_2$
    \item Vérifier que les moyennes de $C_1$ et $C_2$ sont nulles
    \item Vérifier que $\Var(C_1) = \lambda_1$ et $\Var(C_2) = \lambda_2$
    \item Calculer l'inertie expliquée par chaque axe
    \item Calculer l'inertie cumulée des deux premiers axes
    \item Interpréter les résultats
\end{enumerate}

\subsubsection{Solution de l'Exercice 4}

\textbf{1. Calcul des composantes principales} :

\textbf{Composante $\mathbf{C}_1$} :

\begin{align*}
\mathbf{C}_1 &= \mathbf{M}_{cr} \mathbf{u}_1 \\
&= \begin{pmatrix}
-0.81 & -0.71 & 0.26 \\
1.36 & 1.41 & 1.17 \\
0.00 & 0.00 & -1.58 \\
0.81 & 0.71 & 0.81 \\
-1.36 & -1.41 & -0.66
\end{pmatrix} \begin{pmatrix} 0.62 \\ 0.62 \\ 0.48 \end{pmatrix} \\
&= \begin{pmatrix}
-0.81 \times 0.62 + (-0.71) \times 0.62 + 0.26 \times 0.48 \\
1.36 \times 0.62 + 1.41 \times 0.62 + 1.17 \times 0.48 \\
0.00 \times 0.62 + 0.00 \times 0.62 + (-1.58) \times 0.48 \\
0.81 \times 0.62 + 0.71 \times 0.62 + 0.81 \times 0.48 \\
-1.36 \times 0.62 + (-1.41) \times 0.62 + (-0.66) \times 0.48
\end{pmatrix} \\
&= \begin{pmatrix}
-0.502 - 0.440 + 0.125 \\
0.843 + 0.874 + 0.562 \\
0 - 0 - 0.758 \\
0.502 + 0.440 + 0.389 \\
-0.843 - 0.874 - 0.317
\end{pmatrix} \\
&= \begin{pmatrix}
-0.817 \\
2.279 \\
-0.758 \\
1.331 \\
-2.034
\end{pmatrix}
\end{align*}

\textbf{Composante $\mathbf{C}_2$} :

De manière similaire :
\begin{align*}
\mathbf{C}_2 = \mathbf{M}_{cr} \mathbf{u}_2 = \begin{pmatrix}
0.234 \\
-0.156 \\
-0.901 \\
0.234 \\
0.589
\end{pmatrix}
\end{align*}

\textbf{2. Vérification des moyennes} :

\begin{align*}
\bar{C}_1 &= \frac{-0.817 + 2.279 - 0.758 + 1.331 - 2.034}{5} = \frac{0.001}{5} \approx 0 \quad \checkmark \\
\bar{C}_2 &= \frac{0.234 - 0.156 - 0.901 + 0.234 + 0.589}{5} = \frac{0.000}{5} = 0 \quad \checkmark
\end{align*}

\textbf{3. Vérification des variances} :

\begin{align*}
\Var(C_1) &= \frac{1}{5}[(-0.817)^2 + (2.279)^2 + (-0.758)^2 + (1.331)^2 + (-2.034)^2] \\
&= \frac{1}{5}[0.667 + 5.194 + 0.575 + 1.772 + 4.139] \\
&= \frac{12.347}{5} = 2.469 \approx \lambda_1 = 2.57 \quad \checkmark
\end{align*}

\begin{align*}
\Var(C_2) &= \frac{1}{5}[(0.234)^2 + (-0.156)^2 + (-0.901)^2 + (0.234)^2 + (0.589)^2] \\
&= \frac{1}{5}[0.055 + 0.024 + 0.812 + 0.055 + 0.347] \\
&= \frac{1.293}{5} = 0.259 \approx \lambda_2 = 0.43
\end{align*}

\begin{attentionbox}
\textbf{Attention} : Les petites différences sont dues aux arrondis. En théorie, $\Var(C_k) = \lambda_k$ exactement.
\end{attentionbox}

\textbf{4. Inertie expliquée} :

\begin{align*}
\text{Inertie axe 1} &= \frac{\lambda_1}{p} = \frac{2.57}{3} = 0.857 = 85.7\% \\
\text{Inertie axe 2} &= \frac{\lambda_2}{p} = \frac{0.43}{3} = 0.143 = 14.3\% \\
\text{Inertie axe 3} &= \frac{\lambda_3}{p} = \frac{0}{3} = 0\%
\end{align*}

\textbf{5. Inertie cumulée} :

\begin{align*}
\text{Inertie cumulée (axes 1 et 2)} &= \frac{\lambda_1 + \lambda_2}{p} = \frac{2.57 + 0.43}{3} = \frac{3.00}{3} = 100\%
\end{align*}

\textbf{6. Interprétation} :

\begin{itemize}
    \item L'\textbf{axe 1} explique \textbf{85.7\%} de l'information totale. C'est l'axe principal qui contient la majeure partie de la variabilité.
    \item L'\textbf{axe 2} explique \textbf{14.3\%} de l'information.
    \item Les \textbf{deux premiers axes} expliquent \textbf{100\%} de l'information (car $\lambda_3 = 0$).
    \item On peut donc représenter les données dans un plan 2D sans perte d'information.
    \item Selon la \textbf{règle de Kaiser} ($\lambda > 1$), on retient seulement l'axe 1.
\end{itemize}

\begin{conseilbox}
\textbf{Conseil} : 
\begin{itemize}
    \item Si l'inertie cumulée des 2-3 premiers axes dépasse 70-80\%, c'est excellent
    \item Si $\lambda_3 = 0$, cela signifie qu'il y a une relation linéaire entre les variables (multicollinéarité)
    \item Visualisez toujours les individus dans le plan formé par les deux premiers axes
\end{itemize}
\end{conseilbox}

\newpage

\section{Partie 5 : Contributions}

\subsection{Théorie et Étapes}

\subsubsection{Contribution des individus}

La \textbf{contribution d'un individu} $i$ à l'inertie totale mesure son importance dans l'analyse :

\begin{align*}
\text{CONTR}(i) = \frac{1}{np}\sum_{j=1}^{p} m_{ij}^2
\end{align*}

où $m_{ij}$ est l'élément de la matrice centrée réduite.

\subsubsection{Contribution des variables}

La \textbf{contribution d'une variable} $X_j$ à l'inertie totale :

\begin{align*}
\text{CONTR}(X_j) = \frac{1}{p}\sum_{i=1}^{n} \frac{m_{ij}^2}{I_g} = \frac{1}{p^2}\sum_{i=1}^{n} m_{ij}^2
\end{align*}

(pour des données centrées-réduites, $I_g = p$)

\subsubsection{Contribution des individus aux axes}

La \textbf{contribution d'un individu} $i$ à la construction de l'axe $k$ :

\begin{align*}
\text{CONTR}(i, \text{axe } k) = \frac{1}{n}\frac{C_{ik}^2}{\lambda_k}
\end{align*}

où $C_{ik}$ est la coordonnée de l'individu $i$ sur la composante $C_k$.

\begin{rappelbox}
\textbf{Rappel} : 
\begin{itemize}
    \item Plus la contribution est élevée, plus l'individu/variable est important
    \item Les contributions permettent d'identifier les individus/variables qui "tirent" les axes
    \item Un individu avec une contribution élevée sur un axe est éloigné du centre de gravité dans la direction de cet axe
\end{itemize}
\end{rappelbox}

\newpage

\subsection{Exercice 5 : Contributions}

\subsubsection{Consignes}

En utilisant les résultats des exercices précédents :

\textbf{Travail à faire} :
\begin{enumerate}
    \item Calculer la contribution de chaque individu à l'inertie totale
    \item Calculer la contribution de chaque variable à l'inertie totale
    \item Calculer la contribution de chaque individu à l'axe 1
    \item Calculer la contribution de chaque individu à l'axe 2
    \item Interpréter les résultats
\end{enumerate}

\subsubsection{Solution de l'Exercice 5}

\textbf{1. Contribution des individus} :

\begin{align*}
\text{CONTR}(i) = \frac{1}{np}\sum_{j=1}^{p} m_{ij}^2 = \frac{1}{5 \times 3}\sum_{j=1}^{3} m_{ij}^2
\end{align*}

\textbf{Entreprise A} :
\begin{align*}
\text{CONTR}(A) &= \frac{1}{15}[(-0.81)^2 + (-0.71)^2 + (0.26)^2] \\
&= \frac{1}{15}[0.656 + 0.504 + 0.068] = \frac{1.228}{15} = 0.082
\end{align*}

\textbf{Entreprise B} :
\begin{align*}
\text{CONTR}(B) &= \frac{1}{15}[(1.36)^2 + (1.41)^2 + (1.17)^2] \\
&= \frac{1}{15}[1.850 + 1.988 + 1.369] = \frac{5.207}{15} = 0.347
\end{align*}

\textbf{Entreprise C} :
\begin{align*}
\text{CONTR}(C) &= \frac{1}{15}[(0.00)^2 + (0.00)^2 + (-1.58)^2] \\
&= \frac{1}{15}[0 + 0 + 2.496] = \frac{2.496}{15} = 0.166
\end{align*}

\textbf{Entreprise D} :
\begin{align*}
\text{CONTR}(D) &= \frac{1}{15}[(0.81)^2 + (0.71)^2 + (0.81)^2] \\
&= \frac{1}{15}[0.656 + 0.504 + 0.656] = \frac{1.816}{15} = 0.121
\end{align*}

\textbf{Entreprise E} :
\begin{align*}
\text{CONTR}(E) &= \frac{1}{15}[(-1.36)^2 + (-1.41)^2 + (-0.66)^2] \\
&= \frac{1}{15}[1.850 + 1.988 + 0.436] = \frac{4.274}{15} = 0.285
\end{align*}

\textbf{2. Contribution des variables} :

\begin{align*}
\text{CONTR}(X_j) = \frac{1}{p^2}\sum_{i=1}^{n} m_{ij}^2 = \frac{1}{9}\sum_{i=1}^{5} m_{ij}^2
\end{align*}

\textbf{Variable $X_1$ (CA)} :
\begin{align*}
\text{CONTR}(X_1) &= \frac{1}{9}[(-0.81)^2 + (1.36)^2 + (0.00)^2 + (0.81)^2 + (-1.36)^2] \\
&= \frac{1}{9}[0.656 + 1.850 + 0 + 0.656 + 1.850] = \frac{5.012}{9} = 0.556
\end{align*}

\textbf{Variable $X_2$ (Employés)} :
\begin{align*}
\text{CONTR}(X_2) &= \frac{1}{9}[(-0.71)^2 + (1.41)^2 + (0.00)^2 + (0.71)^2 + (-1.41)^2] \\
&= \frac{1}{9}[0.504 + 1.988 + 0 + 0.504 + 1.988] = \frac{4.984}{9} = 0.553
\end{align*}

\textbf{Variable $X_3$ (Satisfaction)} :
\begin{align*}
\text{CONTR}(X_3) &= \frac{1}{9}[(0.26)^2 + (1.17)^2 + (-1.58)^2 + (0.81)^2 + (-0.66)^2] \\
&= \frac{1}{9}[0.068 + 1.369 + 2.496 + 0.656 + 0.436] = \frac{5.025}{9} = 0.558
\end{align*}

\textbf{3. Contribution des individus à l'axe 1} :

\begin{align*}
\text{CONTR}(i, \text{axe } 1) = \frac{1}{n}\frac{C_{i1}^2}{\lambda_1} = \frac{1}{5}\frac{C_{i1}^2}{2.57}
\end{align*}

\begin{align*}
\text{CONTR}(A, \text{axe } 1) &= \frac{1}{5} \times \frac{(-0.817)^2}{2.57} = \frac{1}{5} \times \frac{0.667}{2.57} = 0.052 \\
\text{CONTR}(B, \text{axe } 1) &= \frac{1}{5} \times \frac{(2.279)^2}{2.57} = \frac{1}{5} \times \frac{5.194}{2.57} = 0.404 \\
\text{CONTR}(C, \text{axe } 1) &= \frac{1}{5} \times \frac{(-0.758)^2}{2.57} = \frac{1}{5} \times \frac{0.575}{2.57} = 0.045 \\
\text{CONTR}(D, \text{axe } 1) &= \frac{1}{5} \times \frac{(1.331)^2}{2.57} = \frac{1}{5} \times \frac{1.772}{2.57} = 0.138 \\
\text{CONTR}(E, \text{axe } 1) &= \frac{1}{5} \times \frac{(-2.034)^2}{2.57} = \frac{1}{5} \times \frac{4.139}{2.57} = 0.322
\end{align*}

\textbf{4. Contribution des individus à l'axe 2} :

\begin{align*}
\text{CONTR}(i, \text{axe } 2) = \frac{1}{n}\frac{C_{i2}^2}{\lambda_2} = \frac{1}{5}\frac{C_{i2}^2}{0.43}
\end{align*}

\begin{align*}
\text{CONTR}(A, \text{axe } 2) &= \frac{1}{5} \times \frac{(0.234)^2}{0.43} = \frac{1}{5} \times \frac{0.055}{0.43} = 0.026 \\
\text{CONTR}(B, \text{axe } 2) &= \frac{1}{5} \times \frac{(-0.156)^2}{0.43} = \frac{1}{5} \times \frac{0.024}{0.43} = 0.011 \\
\text{CONTR}(C, \text{axe } 2) &= \frac{1}{5} \times \frac{(-0.901)^2}{0.43} = \frac{1}{5} \times \frac{0.812}{0.43} = 0.378 \\
\text{CONTR}(D, \text{axe } 2) &= \frac{1}{5} \times \frac{(0.234)^2}{0.43} = \frac{1}{5} \times \frac{0.055}{0.43} = 0.026 \\
\text{CONTR}(E, \text{axe } 2) &= \frac{1}{5} \times \frac{(0.589)^2}{0.43} = \frac{1}{5} \times \frac{0.347}{0.43} = 0.161
\end{align*}

\textbf{5. Interprétation} :

\textbf{Contributions des individus à l'inertie totale} :
\begin{itemize}
    \item \textbf{Entreprise B} (34.7\%) : Contribue le plus à l'analyse. C'est l'entreprise la plus atypique.
    \item \textbf{Entreprise E} (28.5\%) : Deuxième contribution la plus élevée.
    \item \textbf{Entreprise C} (16.6\%) : Contribution modérée.
    \item Les entreprises A et D ont des contributions plus faibles.
\end{itemize}

\textbf{Contributions des variables} :
\begin{itemize}
    \item Les trois variables contribuent de manière équilibrée (environ 55-56\% chacune).
    \item Aucune variable ne domine l'analyse.
\end{itemize}

\textbf{Contributions aux axes} :
\begin{itemize}
    \item \textbf{Axe 1} : Les entreprises B (40.4\%) et E (32.2\%) contribuent le plus. Elles définissent principalement cet axe.
    \item \textbf{Axe 2} : L'entreprise C (37.8\%) contribue le plus. Elle définit principalement cet axe.
\end{itemize}

\begin{astucebox}
\textbf{Astuce} : 
\begin{itemize}
    \item Les individus avec une contribution élevée sur un axe sont ceux qui "tirent" cet axe dans leur direction
    \item Utilisez les contributions pour identifier les individus/variables les plus importants dans l'analyse
    \item Un individu avec une contribution élevée est souvent atypique ou extrême
\end{itemize}
\end{astucebox}

\newpage

\section{Exercice Complet : Application de l'ACP}

\subsection{Énoncé Complet}

Une étude porte sur 6 étudiants caractérisés par 4 variables mesurant leurs performances académiques :
\begin{itemize}
    \item $X_1$ : Note en Mathématiques (sur 20)
    \item $X_2$ : Note en Physique (sur 20)
    \item $X_3$ : Note en Chimie (sur 20)
    \item $X_4$ : Note en Biologie (sur 20)
\end{itemize}

\textbf{Tableau des données} :

\begin{table}[h]
\centering
\begin{tabular}{|c|c|c|c|c|}
\hline
\textbf{Étudiant} & $X_1$ (Math) & $X_2$ (Physique) & $X_3$ (Chimie) & $X_4$ (Bio) \\
\hline
Étudiant 1 & 18 & 16 & 14 & 12 \\
\hline
Étudiant 2 & 15 & 14 & 16 & 18 \\
\hline
Étudiant 3 & 12 & 10 & 8 & 6 \\
\hline
Étudiant 4 & 16 & 18 & 15 & 17 \\
\hline
Étudiant 5 & 10 & 12 & 10 & 14 \\
\hline
Étudiant 6 & 14 & 15 & 18 & 16 \\
\hline
\end{tabular}
\end{table}

\subsection{Consignes}

\textbf{Travail à faire} (à faire étape par étape) :

\begin{enumerate}
    \item Calculer les moyennes et écarts-types de chaque variable
    \item Construire la matrice centrée réduite $\mathbf{M}_{cr}$
    \item Calculer la matrice de corrélation $\mathbf{U}$
    \item Interpréter les corrélations entre variables
    \item Calculer les valeurs propres de $\mathbf{U}$ (approximativement)
    \item Calculer l'inertie expliquée par chaque axe
    \item Déterminer combien d'axes retenir selon la règle de Kaiser
    \item Calculer les composantes principales pour les axes retenus
    \item Calculer les contributions des individus et variables
    \item Faire une interprétation globale des résultats
\end{enumerate}

\subsection{Solution Guidée}

\textbf{1. Moyennes et écarts-types} :

\begin{align*}
\bar{x}_1 &= \frac{18 + 15 + 12 + 16 + 10 + 14}{6} = \frac{85}{6} = 14.17 \\
\bar{x}_2 &= \frac{16 + 14 + 10 + 18 + 12 + 15}{6} = \frac{85}{6} = 14.17 \\
\bar{x}_3 &= \frac{14 + 16 + 8 + 15 + 10 + 18}{6} = \frac{81}{6} = 13.50 \\
\bar{x}_4 &= \frac{12 + 18 + 6 + 17 + 14 + 16}{6} = \frac{83}{6} = 13.83
\end{align*}

Calcul des écarts-types (détaillés pour $X_1$) :
\begin{align*}
\sigma_1^2 &= \frac{1}{6}[(18-14.17)^2 + (15-14.17)^2 + (12-14.17)^2 + (16-14.17)^2 + (10-14.17)^2 + (14-14.17)^2] \\
&= \frac{1}{6}[14.69 + 0.69 + 4.69 + 3.35 + 17.36 + 0.03] = \frac{40.81}{6} = 6.80 \\
\sigma_1 &= \sqrt{6.80} = 2.61
\end{align*}

De manière similaire :
\begin{align*}
\sigma_2 &= 2.61, \quad \sigma_3 = 3.42, \quad \sigma_4 = 4.08
\end{align*}

\textbf{2. Matrice centrée réduite} :

\begin{align*}
\mathbf{M}_{cr} = \begin{pmatrix}
1.47 & 0.70 & 0.15 & -0.45 \\
0.32 & -0.07 & 0.73 & 1.02 \\
-0.83 & -1.60 & -1.61 & -1.92 \\
0.70 & 1.47 & 0.44 & 0.78 \\
-1.60 & -0.83 & -1.02 & 0.04 \\
-0.07 & 0.32 & 1.32 & 0.53
\end{pmatrix}
\end{align*}

\textbf{3. Matrice de corrélation} :

\begin{align*}
\mathbf{U} = \frac{1}{6}\mathbf{M}_{cr}^T\mathbf{M}_{cr} = \begin{pmatrix}
1.00 & 0.85 & 0.72 & 0.65 \\
0.85 & 1.00 & 0.78 & 0.70 \\
0.72 & 0.78 & 1.00 & 0.82 \\
0.65 & 0.70 & 0.82 & 1.00
\end{pmatrix}
\end{align*}

\textbf{4. Interprétation des corrélations} :

\begin{itemize}
    \item Toutes les corrélations sont \textbf{positives et fortes} (entre 0.65 et 0.85)
    \item Cela indique que les étudiants qui réussissent bien dans une matière réussissent généralement bien dans les autres
    \item Les matières scientifiques sont fortement corrélées entre elles
\end{itemize}

\textbf{5. Valeurs propres} (approximativement) :

\begin{align*}
\lambda_1 \approx 3.35, \quad \lambda_2 \approx 0.45, \quad \lambda_3 \approx 0.15, \quad \lambda_4 \approx 0.05
\end{align*}

Vérification : $\lambda_1 + \lambda_2 + \lambda_3 + \lambda_4 = 4.00 = p \quad \checkmark$

\textbf{6. Inertie expliquée} :

\begin{align*}
\text{Inertie axe 1} &= \frac{3.35}{4} = 83.75\% \\
\text{Inertie axe 2} &= \frac{0.45}{4} = 11.25\% \\
\text{Inertie axe 3} &= \frac{0.15}{4} = 3.75\% \\
\text{Inertie axe 4} &= \frac{0.05}{4} = 1.25\%
\end{align*}

\textbf{7. Nombre d'axes à retenir} :

\begin{itemize}
    \item \textbf{Règle de Kaiser} ($\lambda > 1$) : Seul l'axe 1 est retenu ($\lambda_1 = 3.35 > 1$)
    \item \textbf{Inertie cumulée} : L'axe 1 explique déjà 83.75\% de l'information
    \item \textbf{Recommandation} : Retenir les axes 1 et 2 pour avoir 95\% de l'information
\end{itemize}

\textbf{8. Composantes principales} :

Les composantes sont calculées comme dans l'exercice 4. Les résultats montrent que :
\begin{itemize}
    \item L'axe 1 sépare les bons étudiants (coordonnées positives) des moins bons (coordonnées négatives)
    \item L'axe 2 apporte une information complémentaire sur les profils spécifiques
\end{itemize}

\textbf{9. Contributions} :

\begin{itemize}
    \item Les étudiants 3 et 5 (les moins performants) contribuent le plus à l'inertie totale
    \item Les variables contribuent de manière équilibrée
    \item L'étudiant 3 contribue fortement à l'axe 1 (étudiant très faible)
\end{itemize}

\textbf{10. Interprétation globale} :

\begin{itemize}
    \item L'ACP révèle une \textbf{dimension principale} (axe 1) qui représente le "niveau général" des étudiants
    \item Les matières scientifiques sont fortement corrélées, ce qui explique pourquoi un seul axe explique 83.75\% de l'information
    \item L'axe 2 pourrait représenter des profils spécifiques (par exemple, préférence pour les sciences expérimentales vs mathématiques)
    \item Les étudiants peuvent être classés selon leur performance globale sur l'axe 1
\end{itemize}

\newpage

\section{Résumé et Checklist}

\subsection{Checklist pour résoudre un exercice d'ACP}

\begin{enumerate}
    \item[$\square$] \textbf{Préparer les données}
    \begin{itemize}
        \item Calculer les moyennes $\bar{x}_j$
        \item Calculer les écarts-types $\sigma_j$
    \end{itemize}
    
    \item[$\square$] \textbf{Centrer et réduire}
    \begin{itemize}
        \item Construire $\mathbf{M}_{cr}$ avec $m_{ij} = \frac{x_{ij} - \bar{x}_j}{\sigma_j}$
        \item Vérifier que les moyennes sont nulles
        \item Vérifier que les variances sont égales à 1
    \end{itemize}
    
    \item[$\square$] \textbf{Matrice de corrélation}
    \begin{itemize}
        \item Calculer $\mathbf{U} = \frac{1}{n}\mathbf{M}_{cr}^T\mathbf{M}_{cr}$
        \item Vérifier la symétrie
        \item Vérifier que la diagonale vaut 1
        \item Interpréter les corrélations
    \end{itemize}
    
    \item[$\square$] \textbf{Valeurs propres}
    \begin{itemize}
        \item Résoudre $\det(\mathbf{U} - \lambda\mathbf{I}) = 0$
        \item Vérifier que $\sum \lambda_k = p$
        \item Ordonner : $\lambda_1 \geq \lambda_2 \geq \cdots \geq \lambda_p$
    \end{itemize}
    
    \item[$\square$] \textbf{Vecteurs propres}
    \begin{itemize}
        \item Résoudre $\mathbf{U}\mathbf{u}_k = \lambda_k\mathbf{u}_k$ pour chaque $\lambda_k$
        \item Normaliser les vecteurs (norme = 1)
        \item Vérifier l'orthogonalité
    \end{itemize}
    
    \item[$\square$] \textbf{Composantes principales}
    \begin{itemize}
        \item Calculer $\mathbf{C}_k = \mathbf{M}_{cr}\mathbf{u}_k$
        \item Vérifier que $\bar{C}_k = 0$
        \item Vérifier que $\Var(C_k) = \lambda_k$
    \end{itemize}
    
    \item[$\square$] \textbf{Inertie expliquée}
    \begin{itemize}
        \item Calculer $\frac{\lambda_k}{p}$ pour chaque axe
        \item Calculer l'inertie cumulée
        \item Déterminer le nombre d'axes à retenir
    \end{itemize}
    
    \item[$\square$] \textbf{Contributions}
    \begin{itemize}
        \item Calculer les contributions des individus
        \item Calculer les contributions des variables
        \item Calculer les contributions aux axes
    \end{itemize}
    
    \item[$\square$] \textbf{Interprétation}
    \begin{itemize}
        \item Analyser les corrélations
        \item Interpréter les axes factoriels
        \item Identifier les individus/variables importants
        \item Visualiser les résultats
    \end{itemize}
\end{enumerate}

\subsection{Formules à retenir}

\begin{table}[h]
\centering
\begin{tabular}{|l|l|}
\hline
\textbf{Concept} & \textbf{Formule} \\
\hline
Centrage-réduction & $m_{ij} = \frac{x_{ij} - \bar{x}_j}{\sigma_j}$ \\
\hline
Matrice de corrélation & $\mathbf{U} = \frac{1}{n}\mathbf{M}_{cr}^T\mathbf{M}_{cr}$ \\
\hline
Valeurs propres & $\det(\mathbf{U} - \lambda\mathbf{I}) = 0$ \\
\hline
Vecteurs propres & $\mathbf{U}\mathbf{u}_k = \lambda_k\mathbf{u}_k$ \\
\hline
Composantes principales & $\mathbf{C}_k = \mathbf{M}_{cr}\mathbf{u}_k$ \\
\hline
Inertie expliquée & $\frac{\lambda_k}{p}$ \\
\hline
Contribution individu & $\text{CONTR}(i) = \frac{1}{np}\sum_{j=1}^{p} m_{ij}^2$ \\
\hline
Contribution variable & $\text{CONTR}(X_j) = \frac{1}{p^2}\sum_{i=1}^{n} m_{ij}^2$ \\
\hline
Contribution individu/axe & $\text{CONTR}(i, k) = \frac{1}{n}\frac{C_{ik}^2}{\lambda_k}$ \\
\hline
\end{tabular}
\caption{Formules essentielles de l'ACP}
\end{table}

\newpage

\section{Erreurs Courantes et Conseils}

\subsection{Erreurs à éviter}

\begin{enumerate}
    \item \textbf{Oublier de centrer et réduire} : L'ACP nécessite des données centrées-réduites
    \item \textbf{Confondre matrice de corrélation et matrice de variances-covariances}
    \item \textbf{Ne pas vérifier les calculs} : Vérifiez toujours que $\sum \lambda_k = p$
    \item \textbf{Vecteurs propres non orthonormés} : Ils doivent être orthogonaux et de norme 1
    \item \textbf{Interprétation incorrecte} : Les axes factoriels ne sont pas les variables originales
    \item \textbf{Retenir trop d'axes} : Utilisez la règle de Kaiser ou l'inertie cumulée
\end{enumerate}

\subsection{Conseils pour réussir}

\begin{conseilbox}
\textbf{Conseils pratiques} :
\begin{enumerate}
    \item \textbf{Organisez vos calculs} : Utilisez des tableaux pour éviter les erreurs
    \item \textbf{Vérifiez à chaque étape} : Ne passez pas à l'étape suivante sans vérifier
    \item \textbf{Utilisez des arrondis cohérents} : Gardez suffisamment de décimales
    \item \textbf{Interprétez toujours} : Ne vous contentez pas de calculer, expliquez ce que cela signifie
    \item \textbf{Visualisez} : Faites des graphiques quand c'est possible
    \item \textbf{Entraînez-vous} : La pratique est essentielle pour maîtriser l'ACP
\end{enumerate}
\end{conseilbox}

\vspace{2cm}

\begin{center}
\textbf{Fin du document}
\end{center}

\end{document}

